% Partie : logique du premier ordre
% Chapitre : complétude
% Section : identité

\documentclass[../../../include/open-logic-section]{subfiles}

\begin{document}

\olfileid[fr]{fol}{com}{ide}
\olsection{Égalité}

\begin{explain}
La construction du modèle des termes donnée à la section précédente suffit
à établir la complétude de la logique du premier ordre pour les ensembles
$\Gamma$ dont les formules ne contiennent pas~$\eq$. Le modèle des termes
satisfait tout !!{sentence}~$!A\in\Gamma^*$ qui ne contient pas~$\eq$
(et donc toute formule~$!A\in\Gamma$). Cependant, cette construction ne
convient plus en présence de~$\eq$. En effet, $\Gamma^*$ peut alors
contenir !!a{sentence}~$\eq[t][t']$, alors que, dans le modèle des
termes, la valeur d'un terme est ce terme lui-même. Si $t$ et $t'$ sont des
termes distincts, leurs valeurs dans ce modèle, à savoir $t$ et $t'$, sont
donc distinctes, et $\eq[t][t']$ est fausse. Nous remédions à cela en
passant au quotient.
\end{explain}

\begin{defn}
  Soit $\Gamma^*$ un ensemble cohérent et !!{complete} d'énoncés du
  langage~$\Lang L$. Nous définissons la relation $\approx$ sur l'ensemble
  des termes clos de~$\Lang L$ par
  \[
  t\approx t' \quad\Longleftrightarrow\quad \eq[t][t']\in\Gamma^*.
  \]
\end{defn}

\begin{prop}
\ollabel{prop:approx-equiv}
La relation $\approx$ possède les propriétés suivantes :
\begin{enumerate}
\item $\approx$ est réflexive.
\item $\approx$ est symétrique.
\item $\approx$ est transitive.
\item Si $t\approx t'$, si $f$ est !!a{function} et si $t_1$, \dots,
  $t_{i-1}$, $t_{i+1}$, \dots, $t_n$ sont des termes clos, alors
\[
\Atom{f}{t_1,\dots,t_{i-1},t,t_{i+1},\dots,t_n}\approx
\Atom{f}{t_1,\dots,t_{i-1},t',t_{i+1},\dots,t_n}.
\]
\item Si $t\approx t'$, si $R$ est !!a{predicate} et si $t_1$, \dots,
  $t_{i-1}$, $t_{i+1}$, \dots, $t_n$ sont des termes clos, alors
\begin{multline*}
\Atom{R}{t_1,\dots,t_{i-1},t,t_{i+1},\dots,t_n}\in\Gamma^*
\quad\Longleftrightarrow\quad\\
\Atom{R}{t_1,\dots,t_{i-1},t',t_{i+1},\dots,t_n}\in\Gamma^*.
\end{multline*}
\end{enumerate}
\end{prop}

\begin{proof}
Puisque $\Gamma^*$ est cohérent et !!{complete},
$\eq[t][t']\in\Gamma^*$ si et seulement si
$\Gamma^*\Proves\eq[t][t']$, par réflexivité et par la
\olref[ccs]{prop:ccs}\olref[ccs]{prop:ccs-prov-in}.
Il suffit donc d'établir les faits suivants. Dans le dernier point,
la réciproque s'obtient en échangeant $t$ et~$t'$ par symétrie.
\begin{enumerate}
\item $\Gamma^*\Proves\eq[t][t]$ pour tout terme clos~$t$.
\item Si $\Gamma^*\Proves\eq[t][t']$, alors
  $\Gamma^*\Proves\eq[t'][t]$.
\item Si $\Gamma^*\Proves\eq[t][t']$ et
  $\Gamma^*\Proves\eq[t'][t'']$, alors
  $\Gamma^*\Proves\eq[t][t'']$.
\item Si $\Gamma^*\Proves\eq[t][t']$, alors
\[
\Gamma^*\Proves
\eq[\Atom{f}{t_1,\dots,t_{i-1},t,t_{i+1},\dots,t_n}]
   [\Atom{f}{t_1,\dots,t_{i-1},t',t_{i+1},\dots,t_n}]
\]
pour tout !!{function} $n$-aire~$f$ et tous termes clos $t_1$, \dots,
$t_{i-1}$, $t_{i+1}$, \dots,~$t_n$.
\item Si $\Gamma^*\Proves\eq[t][t']$ et
$\Gamma^*\Proves\Atom{R}{t_1,\dots,t_{i-1},t,t_{i+1},\dots,t_n}$,
alors
$\Gamma^*\Proves\Atom{R}{t_1,\dots,t_{i-1},t',t_{i+1},\dots,t_n}$
pour tout !!{predicate} $n$-aire~$R$ et tous termes clos $t_1$, \dots,
$t_{i-1}$, $t_{i+1}$, \dots,~$t_n$.
\end{enumerate}
\end{proof}

\begin{prob}
Complétez la démonstration de la~\olref[fol][com][ide]{prop:approx-equiv}.
\end{prob}

\begin{defn}
Supposons que $\Gamma^*$ soit un ensemble cohérent et !!{complete} dans
le langage~$\Lang L$, que $t$ soit un terme clos et que $\approx$ soit
la relation définie ci-dessus. Posons
$T_L=\Setabs{u}{u\in\Trm[L]\text{ et }u\text{ est clos}}$.
Alors :
\[
\equivrep{t}{\approx}=
\Setabs{t'}{t'\in T_L,\ t\approx t'}
\]
et
$\equivclass{T_L}{\approx}=
\Setabs{\equivrep{t}{\approx}}{t\in T_L}$.
\end{defn}

\begin{defn}
\ollabel{defn:term-model-factor}
Soit $\Struct M=\Struct M(\Gamma^*)$ le modèle des termes associé à
$\Gamma^*$ dans la~\olref[mod]{defn:termmodel}. Alors
$\Struct{\equivclass{M}{\approx}}$ est la !!{structure} suivante :
\begin{enumerate}
\item $\Domain{\equivclass{M}{\approx}}=
  \equivclass{T_L}{\approx}$.
\item $\Assign{c}{\equivclass{M}{\approx}}=
  \equivrep{c}{\approx}$.
\item $\Assign{f}{\equivclass{M}{\approx}}
  (\equivrep{t_1}{\approx},\dots,\equivrep{t_n}{\approx})=
  \equivrep{\Atom{f}{t_1,\dots,t_n}}{\approx}$.
\item $\tuple{\equivrep{t_1}{\approx},\dots,\equivrep{t_n}{\approx}}
  \in\Assign{R}{\equivclass{M}{\approx}}$ si et seulement si
  $\Sat{M}{\Atom{R}{t_1,\dots,t_n}}$, c'est-à-dire si et seulement si
  $\Atom{R}{t_1,\dots,t_n}\in\Gamma^*$.
\end{enumerate}
\end{defn}

\begin{explain}
Pour définir $\Assign{f}{\equivclass{M}{\approx}}$ et
$\Assign{R}{\equivclass{M}{\approx}}$ sur les éléments de
$\equivclass{T_L}{\approx}$, nous avons désigné ces éléments sous la
forme $\equivrep{t}{\approx}$, c'est-à-dire au moyen de
\emph{représentants}~$t\in\equivrep{t}{\approx}$. Il faut vérifier
que les définitions ne dépendent pas du choix de ces représentants :
d'autres termes~$t'$ de la même classe
($\equivrep{t}{\approx}=\equivrep{t'}{\approx}$) doivent donner le
même résultat. Par exemple, si $R$ est un !!{predicate} unaire, la
dernière clause dit que
$\equivrep{t}{\approx}\in\Assign{R}{\equivclass{M}{\approx}}$
si et seulement si $\Sat{M}{\Atom{R}{t}}$. S'il existait un autre
terme~$t'$ tel que $t\approx t'$ mais
$\Sat/{M}{\Atom{R}{t'}}$,
la même définition imposerait
$\equivrep{t'}{\approx}\notin
\Assign{R}{\equivclass{M}{\approx}}$.
Or $t\approx t'$ entraîne
$\equivrep{t}{\approx}=\equivrep{t'}{\approx}$, et une même classe
ne peut à la fois appartenir et ne pas appartenir à cette relation.
La \olref{prop:approx-equiv} exclut précisément ce cas.
\end{explain}

\begin{prop}
La !!{structure} $\Struct{\equivclass{M}{\approx}}$ est bien définie.
Plus précisément, si $t_1$, \dots, $t_n$, $t_1'$, \dots, $t_n'$ sont
des termes clos tels que $t_i\approx t_i'$ pour chaque $i=1,\dots,n$,
alors :
\begin{enumerate}
\item $\equivrep{\Atom{f}{t_1,\dots,t_n}}{\approx}=
  \equivrep{\Atom{f}{t_1',\dots,t_n'}}{\approx}$, autrement dit
  \[
  \Atom{f}{t_1,\dots,t_n}\approx
  \Atom{f}{t_1',\dots,t_n'};
  \]
\item $\Sat{M}{\Atom{R}{t_1,\dots,t_n}}$ si et seulement si
  $\Sat{M}{\Atom{R}{t_1',\dots,t_n'}}$, autrement dit
  \[
  \Atom{R}{t_1,\dots,t_n}\in\Gamma^*
  \quad\Longleftrightarrow\quad
  \Atom{R}{t_1',\dots,t_n'}\in\Gamma^*.
  \]
\end{enumerate}
\end{prop}

\begin{proof}
La \olref{prop:approx-equiv} s'applique successivement
aux $n$ coordonnées, en remplaçant chaque $t_i$ par~$t_i'$. La
transitivité de~$\approx$ relie les termes obtenus aux deux extrémités,
et l'équivalence d'appartenance se transmet à chaque remplacement.
\end{proof}

Comme dans le cas du modèle des termes, nous avons besoin du lemme
suivant avant de démontrer le lemme de vérité.

\begin{lem}
\ollabel{lem:val-in-termmodel-factored}
Soit $\Struct M=\Struct M(\Gamma^*)$. Pour tout terme clos~$t$,
$\Value{t}{\equivclass{M}{\approx}}=\equivrep{t}{\approx}$.
\end{lem}
\begin{proof}
La démonstration est analogue à celle du
\olref[mod]{lem:val-in-termmodel}.
\end{proof}

\begin{prob}
Complétez la démonstration du
\olref[fol][com][ide]{lem:val-in-termmodel-factored}.
\end{prob}

\begin{lem}
\ollabel{lem:truth}
Soit $\Gamma^*$ un ensemble cohérent, !!{complete} et saturé
comme dans la définition du modèle des termes.
Pour tout énoncé~$!A$,
$\Sat{\equivclass{M}{\approx}}{!A}$ si et seulement si
$!A\in\Gamma^*$.
\end{lem}

\begin{proof}
Raisonnons par récurrence sur le nombre de connecteurs et de
quantificateurs de~$!A$, comme dans la démonstration du
\olref[mod]{lem:truth}. Pour les atomes relationnels, la définition du
quotient et le lemme sur la valeur des termes donnent l'équivalence.
Les cas quantifiés utilisent les instances closes, puisque chaque
élément du quotient est désigné par l'un de ses représentants.
Reste le cas de l'égalité, où $!A\ident\eq[t][t']$.
\begin{align*}
\Sat{\equivclass{M}{\approx}}{\eq[t][t']}
&\Longleftrightarrow \equivrep{t}{\approx}=\equivrep{t'}{\approx}
  &&\text{(par le \olref{lem:val-in-termmodel-factored})}\\
&\Longleftrightarrow t\approx t'
  &&\text{(par définition des classes d'équivalence)}\\
&\Longleftrightarrow \eq[t][t']\in\Gamma^*
  &&\text{(par définition de $\approx$).}
\end{align*}
\end{proof}

\begin{digress}
Dans le langage enrichi de l'expansion de Henkin, le modèle des
termes $\Struct{M(\Gamma^*)}$ est infini dénombrable. Son quotient
$\Struct{\equivclass{M}{\approx}}$ peut être fini s'il n'existe
qu'un nombre fini de classes~$\equivrep{t}{\approx}$. C'est
prévisible : $\Gamma$ peut contenir des !!{sentence}s qui imposent
que toute !!{structure} les satisfaisant soit finie. Par exemple,
$\lforall[x][\lforall[y][x=y]]$ est un !!{sentence} cohérent.
Il est satisfait exactement dans les !!{structure}s dont le
!!{domain} contient un seul !!{element}.
\end{digress}

\end{document}
